\section{Related work}
\label{sec:related}

Pollard's rho method~\cite{pollard1978} and Shanks' baby-step
giant-step~\cite{shanks1971} are the generic algorithms. Teske's
$r$-adding walk~\cite{teske2001} is the mixing rule this library uses;
van Oorschot--Wiener~\cite{vanoorschot1999} is the parallelisation;
Bernstein--Lange--Schwabe~\cite{bls2011} is the correct use of the
negation map. Two grumpy giants and a baby~\cite{bernsteinlange2012grumpy}
and free precomputation~\cite{bernsteinlange2012small,bernsteinlange2013}
are priced as neighbours, not as replacements. Wiener--Zuccherato and
Duursma--Gaudry--Morain~\cite{wienerzuccherato1998,duursma1999} folded
rho by the automorphism group; GLV~\cite{glv2001} is the constructive
counterpart that this library applies automatically on $j=0$ and
$j=1728$. Cheon~\cite{cheon2006,cheon2010} is a different problem ---
DLP with auxiliary inputs --- and is labelled as such.

The linear sieve~\cite{cos1986} and Lanczos over finite
fields~\cite{lamacchia1990} are the $(\ZZ/p\ZZ)^*$ solver. Semaev
summation polynomials~\cite{semaev2004}, Gaudry's index calculus for
abelian varieties of small dimension~\cite{gaudry2009}, Diem's analysis
of the elliptic discrete logarithm~\cite{diem2011}, and Joux--Vitse
pair decompositions~\cite{jouxvitse2013} are the elliptic side. The
contribution here is not a new algorithm; it is a measurement in one
unit, with every phase priced, against a floor that cannot be tuned
away.

ECC2K-130 is Bailey et al.~\cite{bailey2009}, building on Harley's
ECC2K-95~\cite{harley1998} and Certicom's challenge
list~\cite{certicom1997}. Itoh--Tsujii
inversion~\cite{itoh1988} and Montgomery's simultaneous
inversion~\cite{montgomery1987} are the two inversion tricks the
designs spend most of their area and most of their GPU time on. The
FPGA and GPU numbers in Sections~\ref{sec:ecc2k-gpu} and~\ref{sec:ecc2k-fpga}
are this programme's, cited to their artefacts; they are not a claim
to have broken the challenge.

The reporting rule --- boundary, one unit, four classes, phases priced
separately, extrapolations marked --- is the sibling study library's
standing convention, written to stop a thread from spending weeks
improving a headline constant while the ratio to the floor stays
flat. Wesolowski's supersingular-isogeny result in time
$p^{1/3+o(1)}$ is the target-result profile that convention is aimed
at: exponent-first, heuristics numbered, phases booked as
per-attempt-cost times inverse success probability. This paper is a
measurement paper, not that theorem; the profile is named so the
negative elliptic result is read as a finished thread, not as a
pause.
